\section{Reversible updates and retention interventions}
\label{sec:mechanics}

\subsection{Reflection with a batch tape}

Write the trainable parameters and momenta as \(q,p\in\R^N\).
For a fixed batch \(b\), let \(U_b(q)\) be its mean loss,
\(g_b=\nabla U_b\), and \(V=M^{-1}\) for a constant symmetric
positive-definite mass \(M\). The drift--kick--drift
(\DKD{}) position-Verlet step is
\begin{equation}
 x=q+\tfrac h2 Vp,\qquad p'=p-hg_b(x),\qquad
 q'=x+\tfrac h2 Vp'.
 \label{eq:dkd}
\end{equation}
It uses a single fresh gradient and returns synchronized \(q'\) and \(p'\).
From \((q',-p')\), the first drift reaches \(x\). The same force yields
\(-p\), and the final drift reaches \(q\). Therefore \(S_b^{-1}=R S_b R\)
with \(R(q,p)=(q,-p)\). Inverting a sequence reflects at its endpoints and
applies the original batches in reverse order. Replay needs the exact
realized features, labels, order, and any stochastic choices. Event IDs
suffice only with an immutable or deterministically regenerable tape.

Each real-arithmetic drift or kick has a shear Jacobian with a symmetric
off-diagonal block, respectively proportional to \(V\) or to
\(\nabla^2 U_b\). It preserves the canonical two-form, as does the
composition. Phase~1 uses the symmetric kick--drift--kick variant with
two force sites on the same batch. Its final force can be reused at the
next first kick only if the force law is identical at that point.
Arbitrary batch changes break that reuse. Phases~2 and~3 use
\eqref{eq:dkd} to match one-force baselines.

\subsection{Integer shears and a diagnostic for native error}

For unit mass, encode \(q=2^{-s_q}Q\) and \(p=2^{-s_p}P\). With
round-to-nearest, ties-to-even \(r\), the state updates are
\begin{align}
 X &= Q+r\!\left(\tfrac h2\,2^{s_q-s_p}P\right), \\
 P'&= P+r\!\left(-h\,2^{s_p}g_b(\operatorname{f32}(2^{-s_q}X))\right),\\
 Q'&= X+r\!\left(\tfrac h2\,2^{s_q-s_p}P'\right).
 \label{eq:lattice}
\end{align}
The decoder gives the model a binary32 midpoint, while integers
carry the stored endpoints. Multipliers and rounded increments are computed in
binary64 under the configured bounds. Before any store, both the
increment and the prospective sum must fit the permitted range.
Overflow, wrapping, and state rescaling are disallowed. The starting
state must decode to its original binary32 bits.

On admissible steps, \(r(-a)=-r(a)\) makes each reflected drift undo
the corresponding forward translation exactly. The identical midpoint
and batch give the identical binary32 force, so the reflected
kick also undoes the integer translation. Thus the quantized shears
give bit-exact round trips on the checked domain. The classical
symplectic statement applies to the differentiable real-arithmetic
map in \eqref{eq:dkd}. The lattice proof establishes exact inversion
of the discrete map. Its forward deviation from the real or native
binary32 trajectory requires a separate check.

For native arithmetic, let \(r_i\) be the local rounding injection at
substep \(i\) and \(L_i\) a local Lipschitz bound for the corresponding
real map. Iterating the one-step error bound gives
\begin{equation}
 \|e_m\|\le \|e_0\|\prod_{i=1}^{m}L_i
  +\sum_{i=1}^{m}\|r_i\|\prod_{j=i+1}^{m}L_j .
 \label{eq:error}
\end{equation}
A tolerance crossing alone therefore says little about the form of
growth. For the Transformer, we propagate a fixed seeded phase
direction by the tangent map, evaluating a Hessian--vector product
at the actual midpoint of each step. We renormalize the direction
and accumulate log stretch. A centered finite-difference control
checks that propagation independently. The directional finite-time
rate measures that perturbation along that tape. Uniform Lipschitz
bounds and asymptotic Lyapunov exponents require separate evidence.
In the unit-mass case, the recurrence for perturbation \((u,v)\) is
\[
 a=u+\tfrac h2v,\qquad
 v'=v-h\nabla^2 U_b(x)a,\qquad u'=a+\tfrac h2v' .
\]
We compare the recorded cumulative stretch with native replay
error \(\rho_T=\|(\Delta q,\Delta p)\|_2/
\max(1,\|(q_0,p_0)\|_2)\). The small-neural protocol uses
absolute component errors and a separate correct-count criterion.

\subsection{Fisher-guided interventions}

The empirical Fisher diagonal \(f\) averages squared individual
prior-task score derivatives at the pretrained \(q_*\). The
lower median, including ties, gives diagonal projectors \(P_L\)
and \(P_H=I-P_L\). The Phase~1 split brackets the conservative
update by half steps of \(D_{\tau,G}(q,p)=(q,\exp(-\tau\gamma G)p)\)
with \(G=I\) or \(G=P_L\). At positive \(\gamma\), inverting
the damping requires an expanding exponential. Replay uses
the zero-damping map. A passive scalar reservoir can record
removed kinetic energy, yielding an energy-degenerate metric
bracket and nondecreasing reservoir entropy. The projected brake
leaves the high-Fisher momentum block free to move its corresponding
position coordinates. This is why
subspace selectivity does not itself guarantee retention.

Phase~2 separates orientation, inertia, and restoring force.
An energy-matched coordinator computes requested isotropic
removal \(D_0\) on its own arm and low- and high-subspace
capacities \(K_L,K_H\) on their separate trajectories. All arms
remove \(D=\min(D_0,K_L,K_H)\). For arm projector \(P\) with
\(K_P=\|Pp\|^2/2>0\), it updates momentum by
\[
  p^+=(I-P)p+\sqrt{1-D/K_P}\,Pp .
\]
Zero capacity and removal give the identity. The cap preserves
equal realized removal, while changing the nominal isotropic
strength whenever a projected arm lacks capacity.

For the Fisher-mass intervention, let \(W\) be the frozen nonnegative
Fisher diagonal normalized by its mean with a numerical floor.
Use \(M=I+\alpha W\) in \eqref{eq:dkd}. For the GENERIC
intervention, add an anchor potential
\(A(q)=(q-q_*)^\top W(q-q_*)/2\) and passive energy and entropy
coordinates \(e,\sigma\). In variables
\((q,p,e,\sigma-A(q))\), the Poisson part is canonical in \(q,p\)
and leaves the other two coordinates fixed. The symmetric
positive-semidefinite metric part annihilates
\(\nabla(U_b+p^\top Vp/2+e)\). It induces
\(\dot q=-\lambda g W(q-q_*)\), fixes \(p,\sigma\), and
balances the change in \(U_b\) through \(e\). For duration \(\tau\)
the anchor solution is
\[
 q^+=q_*+\exp(-\tau\lambda g W)(q-q_*),\qquad
 e^+=e+U_b(q)-U_b(q^+) .
\]
Two exact half-anchor flows bracket each Fisher-mass DKD step.
At gate \(g=0\) this reduces to its conservative mass control.
The bracket identities concern the continuous flows for one batch.
The discrete conservative step has its usual energy error, and
changing batches can supply external work.
